\subsection{Observed fibroblasts define the material available for temporal prediction}

Three healthy volunteers supplied intact skin and acute-wound biopsies at days 1, 7 and 30, yielding 12 specimens and 12,259 author-annotated fibroblasts (Figure~\ref{fig:p2f1}A\label{call:p2f1}). The computational route in A retains the 7,002-gene input dimension despite coverage of only 5,383 genes, then distinguishes the frozen mean-coordinate encoder, training-only scaling, product-coupled field fitting and numerical prediction. Its globally held-out day-7 task does not depict biological cell paths or the separate donor-transfer protocol. A common projection displays these observed cells by collection time (Figure~\ref{fig:p2f1}B) and by donor (Figure~\ref{fig:p2f1}C). Its first two principal components explain 43.5\% and 10.6\% of the variance in the training set that excludes day 7. Thus, nearly half of the training variance remains outside the displayed plane; visual overlap or separation cannot replace evaluation in all 15 coordinates.
\paperfigure{figure1_workflow}{fig:p2f1}{Acute-wound sampling and frozen-coordinate flow. (A) Conceptual repair, not measured histology or closure time, and repeated biopsies of three healthy donors at days 0, 1, 7 and 30. The lower route starts with 7,002-gene proportions: 5,383 panel genes are covered and the remainder zero-filled. A frozen 15-dimensional encoder mean μ is scaled using training cells only. A shared conditional-flow-matching (CFM) field learns from independently sampled product-coupled endpoints; 50-step RK4 predicts the day-7 marginal in scaled μ coordinates, not simplex weights. All day-7 cells are held from fitting and scaling. This is global time holdout, not the separate held-donor task that permits training-donor day-7 cells. (B) A uniform 2,000-cell display subset of 12,259 fibroblasts by time. PCA and scaling exclude day 7; the axes explain 43.5\% and 10.6\% of training variance. (C) The same cells by donor. (D) Twelve measured genes per donor--time biopsy, with log-transformed pseudobulk counts per 10,000 scaled separately by gene; colors do not compare absolute levels between genes. (E) Fibroblast counts per biopsy; cells are not independent donors. (F) Thirty-six seed-0 numerical paths from observed day-1 cells to day-7 rank time over the day-1/day-7 display. Arrows are not tracked cell fates; errors use all 15 coordinates.}


Matched raw-count profiles show the expression of matrix, inflammatory and contractile genes in each donor's four biopsies (Figure~\ref{fig:p2f1}D). The gene-wise color scale compares a gene across samples and does not compare absolute expression between genes. Retained fibroblast counts differ across donor--time combinations (Figure~\ref{fig:p2f1}E), so neither point density nor the number of cells provides independent temporal replication. Projected paths connect observed day-1 starting coordinates to model-predicted coordinates at the prespecified day-7 time (Figure~\ref{fig:p2f1}F). These are numerical field trajectories; the data do not track the corresponding cells through the wound. Sampling counts are documented in Table~\ref{tab:p2biologicalcounts}.

% Allow the next Results subsection to fill the text page while Figure 1
% remains on the immediately following float page (checked by build.py).
\subsection{A held-out acute human time point can be reconstructed under the fixed protocol}

Global timepoint holdout removes day 7 from every donor, whereas donor transfer removes one donor and retains all four times in the others. Both use the frozen representation and the training-only scaling in Equation~\eqref{eq:p2scale}. The product bridge in Equation~\eqref{eq:p2bridge}, field objective in Equation~\eqref{eq:p2cfm} and initial-value prediction in Equation~\eqref{eq:p2ode} specify a reproducible estimator, not an observed cell transition. This distinction determines whether the target time itself has been observed during training; the permitted source, target exclusions and fitting settings are recorded in Table~\ref{tab:p2protocols}.

A fixed-seed refit assessed numerical resolution of the RK4 update in Equation~\eqref{eq:p2rk4} and sensitivity to evaluation sampling. Across 25, 50, 100 and 200 RK4 steps, the target energy distance remained 0.705395 and the improvement remained 49.863\%; the largest discrepancy from the 200-step coordinates was below $6\times10^{-7}$ in root-mean-square units (Table~\ref{tab:p2numerics}). Repeated evaluation at 250, 500 and 1,000 cells per distribution gave median improvements of 48.46\%, 48.39\% and 49.26\%, respectively. Gain exceeded the same-count split-half reference in all 50 draws at each count (Table~\ref{tab:p2sampling}). These results assess one fitted field and the observed cells; the ranges across draws quantify evaluation sensitivity, not additional donor replication or confidence in clinical generalization.

Using the empirical weighting rule in Equation~\eqref{eq:p2weighted} with all cells and the same fixed 50-step prediction, equal cell weights gave a 49.40\% improvement and equal donor mass gave 49.69\%. The three donor-specific improvements under this pooled field were 47.15\%, 46.92\% and 40.07\% (Table~\ref{tab:p2weights}). This sensitivity changes the evaluation measure without refitting the field; it is distinct from the donor-transfer experiment and from the pooled capped-distance estimate.

\FloatBarrier
The result is deliberately bounded. The donors are healthy volunteers, the wounds are acute experimental wounds and the target is one observed time point. The analysis does not contain a DFU patient trajectory or a time-to-closure endpoint. ``Transfer'' here means transfer across the three specified acute-wound donors under a held-out protocol; it does not mean clinical generalization.
Using the empirical version (Equation~\eqref{eq:p2edv}) of the energy discrepancy defined in Equation~\eqref{eq:p2ed}, the finalized held-out Wound7 analysis improved on the stay-still baseline by +49.9\% under Equation~\eqref{eq:p2improvement} (Figure~\ref{fig:p2f2}A\label{call:p2f2}). Dropping one donor at a time preserved improvement in 3/3 retained-donor sensitivity runs; these runs predict the retained donors, not the omitted donor. The held-out donor transfer analysis also passed in 3/3 folds, with a mean improvement of +75.9\% (Figure~\ref{fig:p2f2}B). In the geometry audit, the median same-donor temporal distance was 1.967 and the median same-time donor distance was 0.233, giving a time-to-donor ratio of 8.46 (descriptive 95\% resampling range 6.02--13.11; Figure~\ref{fig:p2f2}C). The mean displacement cosine was +0.948, with a descriptive resampling range of 0.943--0.954. Shared donors make these pair distances dependent, so the ranges do not have established donor-population coverage. Improvement over unchanged source alone does not establish incremental value over other permitted observations. The probability-conservation identity in Equation~\eqref{eq:p2continuity} concerns the constructed training paths, not identification of a biological trajectory. The observed reconstruction gains remain conditional on this representation, interval and three donors (Tables~\ref{tab:p2folds} and~\ref{tab:p2geometrytable}).

\paperfigure{figure2_wound7_transfer}{fig:p2f2}{Global Wound7 reconstruction, donor transfer and donor--time geometry. (A) Bars show predicted and stay-still energy distances to pooled Wound7 after excluding that time globally; relative improvement is 100 times (stay-still minus predicted distance) divided by stay-still distance, giving 49.9\%. (B) Each bar is a genuinely held-out donor, with all four times retained in the other two donors; the mean relative improvement is 75.9\%. Bars are point estimates without confidence intervals. (C) Points show 12 between-donor same-time distances and 18 within-donor between-time distances. Boxes span the interquartile range, centre lines are medians, and whiskers extend to extreme points within 1.5 interquartile ranges. The temporal/donor median ratio is 8.46 (10,000-draw descriptive pair-distance bootstrap 95\% interval 6.02--13.11). All distances use training-standardized latent coordinates. Shared donors make these pairs dependent; three healthy donors do not provide DFU prediction validation.}


\subsection{Simple interpolation is competitive when an entire time point is missing}

A retrospective extension tested whether improving over unchanged source requires a neural field. The additional comparisons hold out either day 7 or day 1 globally, keeping the frozen encoder, rank clock, training-only scaling and optimization budget fixed. For each task, the observed time immediately before the target supplies the source and the next observed time supplies an anchor. The source and anchor belong to the same donor; the held-time cells enter only the evaluation. Error is calculated using every available cell in all 15 coordinates within each donor, then averaged with equal donor weight. The donor-averaged error in Equation~\eqref{eq:p2macroerror} differs from the earlier pooled and capped estimands. The two endpoint controls in Equation~\eqref{eq:p2interpolation} therefore address model complexity under this separate, explicit scoring rule.

For withheld day 7, unchanged source had mean donor error 1.4725. Across three seeds, shared flow reduced it to 0.8110, while within-donor centroid translation reduced it further to 0.6674. An independent-endpoint linear bridge gave 0.9284 (Figure~\ref{fig:p2f3}A\label{call:p2f3}). The corresponding improvements were 44.9\%, 54.7\% and 36.9\%. Withholding day 1 instead trained on intact skin, day 7 and day 30 and predicted from intact skin. The unchanged-source error was 1.3460; shared flow, centroid translation and the independent bridge gave 1.0571, 0.9293 and 1.0959, corresponding to improvements of 21.5\%, 31.0\% and 18.6\% (Figure~\ref{fig:p2f3}B). The complete seed-level errors appear in Table~\ref{tab:p2interpolation}.

\paperfigure{figure3_interpolation_controls}{fig:p2f3}{Simple baselines for globally missing times. All methods use the same frozen representation and training exclusions; error averages full-cell empirical energy distances within the three donors. (A) Withheld day 7, predicted from day 1 with day 30 as the observed anchor. (B) Withheld day 1, predicted from intact skin with day 7 as anchor. Points show the three-seed mean error; horizontal ranges show seed minima and maxima, not confidence intervals. Both tasks use the fixed rank-time interpolation fraction one half. (C) Donor-specific seed-averaged field error minus centroid-translation error; positive values favor translation. The later anchor belongs to the same donor, so this is retrospective interpolation rather than source-only donor transfer. (D) Each field seed's improvement over the corresponding unchanged source. Seeds and repeated distance evaluations provide no additional biological replication. (E) Seed- and donor-averaged energy distance after separately centering each distribution; this diagnoses shape without decomposing total error. (F) Mean and full range of flow-minus-comparator errors across 50 paired draws of 200 source and 200 target cells per donor; positive values favor the comparator. Location--scale uses only endpoint means and coordinate standard deviations. Matched draw ranges are not confidence intervals.}


The centroid translation had lower seed-averaged error in all three donors for withheld day 7 and in two of three for withheld day 1 (Figure~\ref{fig:p2f3}C). The flow error ranged from 0.8044 to 0.8237 across day-7 fits, and from 0.9484 to 1.1231 across day-1 fits (Table~\ref{tab:p2interpolation}). Every fit improved over its unchanged source; Figure~\ref{fig:p2f3}D displays these seed-specific relative improvements. These ranges characterize training variation in the existing cohort; they do not provide six independent biological replications. All donor-specific comparisons are retained in Table~\ref{tab:p2interpolationdonors}.

The simple predictor uses an observed later biopsy from the donor being reconstructed, which is permitted by the global timepoint-holdout task. It would not be available in prospective source-only prediction for a new donor. Conversely, the donor-transfer benchmark observes the target time in the training donors and asks a different question. Its method comparisons cannot be generalized to a globally unobserved time. The added comparison supports temporal information in the measured distributions, while providing no evidence that a neural field is necessary for the two interpolation tasks tested here.

The location--scale predictor in Equation~\eqref{eq:p2locationscale}, which interpolates endpoint standard deviations as well as means, gave errors of 0.6982 for day 7 and 0.9034 for day 1, also lower than the shared-flow errors in both tasks (Figure~\ref{fig:p2f3}A--B). Mean translation was the better of these two deterministic controls for day 7, whereas location--scale interpolation was better for day 1. This ordering is reported without selecting a new predictor for either target.

The distribution-shape diagnostics explain why a single aggregate error does not identify a uniform advantage. For day 7, flow had slightly lower centered energy distance than translation (0.1441 versus 0.1495), despite higher total error. For day 1, flow had higher centered error (0.1220 versus 0.0647), while the independent bridge also had higher centered error (0.1445; Figure~\ref{fig:p2f3}E). The independent bridge's predicted-to-target variance ratios were 0.319 and 0.579 at days 7 and 1, indicating substantial contraction in this comparison (Table~\ref{tab:p2shapediagnostics}). These are separate geometry summaries, not a causal or additive error decomposition.

Matched 200-cell evaluation preserved the lower error of both deterministic endpoint controls in all 50 draws for each held time. Flow-minus-translation error ranged from 0.1132 to 0.1716 for day 7 and 0.0977 to 0.1545 for day 1; flow-minus-location--scale ranges were 0.0847--0.1405 and 0.1194--0.1763. Flow remained better than the independent bridge in every draw (Figure~\ref{fig:p2f3}F; Table~\ref{tab:p2matchedsampling}). The full-cell calculations reproduced all previous donor/seed errors to numerical precision. These results establish robustness of the observed ordering to the tested evaluation samples, conditional on three donors and the frozen fits.

\subsection{Time parameterization changes the held-out result}

The historical single-seed, nearest-grid comparison gave linear days an improvement of −18.2\% relative to unchanged source (Table~\ref{tab:p2axes}). The revised paired audit instead evaluates 12 fields at the exact target time with identical random streams and evaluation rows within each seed. The day-based transforms in Equation~\eqref{eq:p2times} define separate fitted bridges, rather than transforming one fixed field by the chain rule in Equation~\eqref{eq:p2clock}. Mean donor energy distances were 0.8104, 2.1542, 0.8746 and 0.7725 for rank, linear days, square-root days and log days, respectively, versus an unchanged-source distance of 1.4746. The corresponding centroid translations gave 0.6667, 0.9775, 0.7476 and 0.6654, lower than the shared field under every clock on the donor-and-seed average (Table~\ref{tab:p2pairedclock}; Figure~\ref{fig:p2f4}\label{call:p2f4}). These equal-donor, exact-time values have a different estimand and sampling cap from the historical pooled scan. Changing the coordinate changes bridges and optimization targets, so the comparison does not isolate interval compression or select a universal biological clock.

\paperfigure{figure4_time_axes}{fig:p2f4}{Paired time-coordinate sensitivity under global day-7 holdout in three healthy donors. Twelve fields combine four fixed clocks with seeds 0, 1 and 2, each using 8,000 training steps and batch size 256. Network initialization, endpoint draws and interpolation random draws are paired across clocks within seed. All day-7 cells are excluded from fitting and scaling. Energy distance is evaluated at the exact transformed day-7 time after 50 RK4 steps, using the same source and target rows across clocks and seeds and a 1,200-cell cap per donor. Each point is the equal mean of three donor-specific distances; shapes identify training seeds. Thin lines join the same seed, and thick lines show three-seed means. Colors distinguish shared flow, within-donor centroid translation and unchanged source. The centroid control uses the permitted day-1 and day-30 biopsies, not day-7 expression. Constant baseline points overlap. Connecting categorical clock choices does not depict a biological trajectory, and seed spread is not a biological confidence interval.}




\subsection{The training-target distribution challenges the held-donor method ranking}

The historical held-donor benchmark gave shared flow matching a mean energy distance of 0.3947 (approximately 0.395) and gains above the split-half reference in 3/3 folds. Mean displacement, nearest-donor displacement, entropic transport and the conditioned field also met that descriptive criterion. An added predictor instead uses only the training donors' Wound7 distributions, assigning equal donor mass and matching the held-source particle count. Its mean error was 0.1818, with a full range of 0.1479--0.2150 across 50 prediction-sampling draws (Figure~\ref{fig:p2f5}A\label{call:p2f5}). The same historical target indices and scaling were used, so improvement over unchanged source does not establish incremental value from the excluded donor's source expression (Tables~\ref{tab:p2methods}, \ref{tab:p2methodsbydonor} and~\ref{tab:p2targetmarginal}).

The donor-specific results prevent a universal ranking: target-marginal mean errors were 0.1569, 0.1416 and 0.2469 for PWH26, PWH27 and PWH28, whereas shared-field errors were 0.7110, 0.2881 and 0.1850. The field therefore remained better in PWH28, including all 50 matched marginal draws. Those draws reuse the same three donors and quantify prediction-sampling sensitivity rather than population uncertainty.

The context model in Equation~\eqref{eq:p2context} improved all 6/6 training-donor comparisons (Figure~\ref{fig:p2f5}B; Table~\ref{tab:p2insample}) while giving a higher held-out mean than the shared field. Its greater fit to known contexts therefore supplies no evidence of a stable generalization advantage.

\paperfigure{figure5_methods_benchmark}{fig:p2f5}{Method comparison across three held-out donors. (A) Bars are equally donor-weighted mean energy distances for unchanged source, five historical predictors and the added matched-particle target marginal. CFM means conditional flow matching; OT is the entropic transport displacement extension. The target marginal assigns equal mass to the two training donors' Wound7 distributions, uses no held-source expression, and averages 50 draws under the original target samples and scaling. Its mean error is 0.1818 versus 0.3947 for historical shared CFM, but PWH28 favors the field. All historical predictive methods exceed the single split-half gain reference in 3/3 folds; this descriptive criterion is not a significance test or population ranking. (B) Lines connect shared and conditioned errors for six training-donor comparisons (two retained donors in each of three fits). Conditioning improves all six. These comparisons share donors and do not measure unseen-donor transfer. Neither panel shows uncertainty intervals.}


The original neural and transport method comparison used one training seed and successive endpoint random draws. The entropic plan in Equation~\eqref{eq:p2ot} is a separate comparator, not the coupling used to train the flow. The added target-marginal result is a retrospective baseline audit, and the historical fields were not saved as per-cell predictions for new paired target sampling. Neither result selects a model class for DFU patients; independent donors and matched repeated fits are required for that question.

\subsection{Donor-count comparisons require a common evaluation scale}

The initial count curve changed its standardization when the training subset changed, so subtracting its absolute errors mixed information gain with a change of distance geometry. A controlled rerun retained subset-only training normalization but scored inverse-transformed predictions in one common reference per held donor. Across 27 fits, the hierarchical average in Equation~\eqref{eq:p2curve} gave shared-field mean error of 0.5466 with one training donor and 0.3983 with two. Figure~\ref{fig:p2f6}A\label{call:p2f6} shows all three seeds and the seed-0 line (0.5315 to 0.3854). Seed-specific reductions were 0.1462, 0.1558 and 0.1430 (Figure~\ref{fig:p2f6}B). The mean reduction was 0.1483, with all 3/3 donor blocks improving and a descriptive three-block percentile range of {[}0.0732, 0.2727{]} (Table~\ref{tab:p2curve}).

\paperfigure{figure6_donor_curve_stability}{fig:p2f6}{Donor-count comparison in common scoring coordinates. (A) k is the number of training donors. The line joins seed-0 mean held-out errors, 0.5315 at k=1 and 0.3854 at k=2; symbols identify seeds 0, 1 and 2. Each model uses only its selected training donors for fitting and normalization. Predictions return to encoder coordinates before scoring under one reference fitted from both non-held donors, fixed across k and seed and excluded from single-donor fitting. At k=1, two training choices are averaged within held donor before equally weighting the three donors. (B) Points are seed-specific reductions, k=1 minus k=2. The line is the three-seed mean, 0.1483. Shading is a descriptive percentile 95\% range [0.0732, 0.2727] from 4,000 resamples of three donor blocks after seed averaging; it is not an interval for each point. Source and target indices are fixed. The largest configuration-specific seed range is 0.3377. Overlapping training sets limit interval interpretation, and two k values cannot identify a plateau or required cohort size.}


The same common geometry also showed lower error for mean displacement (0.5095 to 0.4186), the full equal-donor target marginal (0.2222 to 0.1676), and its matched-particle version (0.2332 to 0.1818). Unchanged-source error remained 1.4760 at both training sizes, checking that the reference metric did not change with k (Table~\ref{tab:p2seedcurve}). The full target-marginal decrease, however, is exactly the mixture effect in Equation~\eqref{eq:p2mixture}; it cannot independently identify target-relevant information from adding a donor. The neural and mean-displacement comparisons involve different fitted predictors and are not algebraically identical to this target-marginal mixture.

The saved two-donor fields also allow the starting distribution to be changed without changing training, target cells or scoring scale. Averaged across seeds and donors, the actual held-source prediction had error 0.3983. Replacing its input with equal-donor draws of training Wound1 cells reduced mean error to 0.2450; the matched training Wound7 marginal, which does not use any source expression from the held donor, had error 0.1818 (Table~\ref{tab:p2sourceswap}). The donor effects were heterogeneous: real-source versus substituted-source errors were 0.7586 versus 0.2802 for PWH26, 0.2368 versus 0.2732 for PWH27, and 0.1994 versus 0.1816 for PWH28. Thus, PWH27 favored its own starting cells on average, while the other two did not. The substitution is a sensitivity analysis of the same fitted fields, not a comparison of nine independent donors; the 50 cell draws likewise supply no new biological replication. It does not demonstrate a consistent incremental value for the excluded donor's initial expression in this cohort, and it does not imply that such information is universally unhelpful.

The exact target-marginal reduction was 0.054601 across donors and matched one quarter of the training-target mutual discrepancy, with maximum numerical residual below $2.45\times10^{-15}$ (Table~\ref{tab:p2mixtureidentity}). The additional 200-particle source audit gave macro errors of 0.4036 for real sources, 0.3075 when the two individual substituted-source errors were averaged, and 0.2527 for their sampled mixture. PWH28 exposes the distinction: real-source error was 0.2043, lower than either individual substitution (0.2503 and 0.2636), yet the mixed substitution was lower still at 0.1904. PWH27 also favored its real source over both individual alternatives, while PWH26 did not (Table~\ref{tab:p2mixtureaudit}). Thus, the earlier mixed-source comparison combines source replacement with a change in the predicted mixture and cannot isolate individual source-state value. The exact identity does not decompose the gain from a jointly fitted two-donor field. All results remain conditional on the same three donors, fields, target cells and reference geometry.

Within a fixed donor and training subset, the mean training SD was 0.0701, or 3.07 times the mean split-half reference of 0.0228; the largest seed range was 0.3377 (Table~\ref{tab:p2seedconfigs}). The aggregate direction persisted across seeds, while individual fits remained variable. The three-block range is descriptive because the held-donor folds share training donors. Two training sizes cannot determine a plateau or a required cohort size, and seeds, subsets and repeated marginal draws add no independent biological replication.

\subsection{Observable expression errors do not establish a flow advantage over target marginals}

Raw-expression alignment and frozen re-encoding passed for all 12,259 cells, with maximum encoder-mean discrepancy $1.91\times10^{-6}$. On the exact saved targets and all 5,383 common-panel genes, shared-flow total variation was 0.2934 with one training donor and 0.2823 with two, versus 0.3529 for decoded persistence. The decoded mean-displacement control gave 0.2882 and 0.2825; decoded training-target marginals gave 0.2590 and 0.2546. Directly observed training-target marginals were lower still at 0.1279 and 0.1115, without using the excluded donor's source expression (Table~\ref{tab:p2observable}). Flow improved over decoded persistence in all three donor blocks but over decoded target marginals in only one of three, and over observed target marginals in none. At two training donors its total variation was close to mean displacement, while its Jensen--Shannon divergence was higher, so these metrics do not establish a consistent neural advantage. The all-day-7 sensitivity retained the same aggregate interpretation. The observed-target decoder reference had mean total variation 0.2430 against actual expression (Table~\ref{tab:p2decoderreference}); because it uses the target itself, this is neither prediction performance nor an error floor. The results add a mean-expression evaluation within the existing cohort, not independent biological validation or evidence about unassayed genes.

\subsection{Three mouse arms expose model-arm heterogeneity}

The all-arm audit gave different results under the same protocol, using the endpoint-distance diagnostic in Equation~\eqref{eq:p2geometry}. NDB changed by −1.5\% relative to stay-still and did not beat the baseline. PDB improved by +32.1\%, but its POD7 target failed an endpoint-distance ordering check, warning against a simple interpolation interpretation without proving an off-manifold trajectory. GDB improved by +50.6\% with a POD7 target satisfying that ordering check. In aggregate, 2/3 arms beat stay-still, 2/3 met the noise criterion and 2/3 satisfied the ordering check (Figure~\ref{fig:p2f7}A\label{call:p2f7}). Each arm had one animal per time point, and NDB covered 59.7\% of the frozen panel. These arm-level computations are useful for protocol auditing, but they do not establish a mouse biological trajectory, a cross-species equivalence or DFU transfer (Tables~\ref{tab:p2mouse} and~\ref{tab:p2mousecounts}).

The paired baseline and predicted distances show the absolute error scale within each mouse arm (Figure~\ref{fig:p2f7}B); relative improvement alone would hide these differences.

\paperfigure{figure7_mouse_arms}{fig:p2f7}{Mouse POD7 computational audit for the source-labelled NDB, PDB and GDB arms. POD7 is excluded from fitting and scaling; rank-time fields trained at POD0, POD2 and POD30 predict from the observed POD2 source. (A) Relative energy-distance improvements over stay-still are −1.5\%, +32.1\% and +50.6\%. Hatching marks PDB, whose POD7 target fails the endpoint-distance ordering diagnostic. (B) Within each arm, the left grey bar is stay-still and the right coloured bar is prediction; hatching carries the same diagnostic meaning. The screen compares endpoint distances and does not prove or disprove geodesic interpolation. Each arm has one animal per time point and 4,182/7,002 mapped panel genes. Bars have no uncertainty intervals; cell counts do not supply independent animal replication or cross-species clinical validation.}

